% Insert after the description of the LLM metadata schema in the Methodology.

\subsection{Post-extraction taxonomy normalisation}

The LLM-generated fields \texttt{thematique} and
\texttt{nature\_initiative} may contain lexical variants that refer to closely
related concepts. A deterministic post-processing step therefore adds the fields
\texttt{thematique\_normalisee} and
\texttt{nature\_initiative\_normalisee}. The original LLM values are preserved
unchanged, while the new fields map them to a documented French controlled
taxonomy used for filtering and descriptive statistics.

The thematic taxonomy contains fifteen substantive categories, including
\textit{autonomie et avancée en âge}, \textit{santé et prévention},
\textit{lien social et lutte contre l'isolement}, \textit{mobilité et
accessibilité}, \textit{habitat et logement}, and \textit{participation
citoyenne et gouvernance}. The initiative-type taxonomy includes categories such
as \textit{atelier}, \textit{formation}, \textit{événement}, \textit{outil
numérique}, \textit{dispositif ou programme}, and \textit{projet
d'aménagement}. A record may have multiple normalised themes but only one
normalised initiative type.

The normalisation is rule-based and does not call the LLM again. It operates on
the extracted labels rather than on the full source text, which keeps the
transformation reproducible and inspectable. Values that cannot be mapped
unambiguously are assigned to \textit{autre / à revoir} and reported in a
dedicated review file. This conservative choice avoids hiding uncertainty behind
an overly broad automated category while making the taxonomy progressively
improvable through documented manual review.
